\documentclass[english,hebrew]{article}
\usepackage[HE8,T1]{fontenc}
\usepackage[utf8]{inputenc}
\setcounter{secnumdepth}{-2}
\setlength{\parindent}{0bp}
\usepackage{amsmath}
\usepackage{amsthm}
\usepackage{amssymb}

\makeatletter
\theoremstyle{plain}
\newtheorem{thm}{\protect\theoremname}
\theoremstyle{definition}
\newtheorem{problem}[thm]{\protect\problemname}
\theoremstyle{definition}
\newtheorem*{sol*}{\protect\solutionname}

\@ifundefined{date}{}{\date{}}
\usepackage{graphicx}
\usepackage{geometry}
\newcommand{\limi}{\underset{n\rightarrow\infty}{\lim}}
\newcommand{\limxz}{\underset{x\rightarrow x_{0}}{\lim}}
\newcommand{\limfxz}{\underset{x\rightarrow x_{0}}{\lim}f(x)}
\newcommand{\limfxm}{\underset{x\rightarrow x_{0}^{-}}{\lim}f(x)}
\newcommand{\limfxp}{\underset{x\rightarrow x_{0}^{+}}{\lim}f(x)}
\newcommand{\limxm}{\underset{x\rightarrow x_{0}^{-}}{\lim}}
\newcommand{\limxp}{\underset{x\rightarrow x_{0}^{+}}{\lim}}
\newcommand{\sqrm}{M_{m\times m}(\mathbb{F})}
\newcommand{\seqa}{(a_{n})_{n=1}^{\infty}}
\newcommand{\seqx}{(x_{n})_{n=1}^{\infty}}
\newcommand{\funcdr}{f:D\rightarrow\mathbb{R}}
\newcommand{\der}{\limxz\frac{f(x)-f(x_{0})}{x-x_{0}}}
\usepackage[all]{pdfprivacy}

\makeatother

\usepackage{babel}
\addto\captionsenglish{\renewcommand{\problemname}{Problem}}
\addto\captionsenglish{\renewcommand{\solutionname}{Solution}}
\addto\captionsenglish{\renewcommand{\theoremname}{Theorem}}
\addto\captionshebrew{\renewcommand{\problemname}{בעיה}}
\addto\captionshebrew{\renewcommand{\solutionname}{פתרון}}
\addto\captionshebrew{\renewcommand{\theoremname}{משפט}}
\providecommand{\problemname}{בעיה}
\providecommand{\solutionname}{פתרון}
\providecommand{\theoremname}{משפט}

\begin{document}
\newgeometry{top=1cm, bottom=2cm, left=1cm, right=1cm}
\title{שאלות לקראת בוחן {\beginL 2\endL}}
\maketitle
\begin{abstract}
נושאים בלקט: פולינום טיילור עם משפטי השארית וכלל לופיטל. מבחינת חומר
לבוחן, זה לא מכסה את הכל.
\end{abstract}
\newpage{}
\begin{problem}
מצאו את הערך המקסימלי של \L{$n\in\mathbb{N}\cup\{0\}$} שעבורו הגבול
\L{%
\[
\lim_{x\to0}\frac{\ln(1+x^{2})-x\sin(x)}{x^{n}}
\]
}%
קיים במובן הצר, וחשבו את ערכו עבור אותו \L{$n$}.
\end{problem}

\begin{sol*}
נחשב את פולינומי טיילור של הפונקציות במונה\L{%
\[
P_{n,\ln\left(1+x\right),0}(t)=\underset{k=0}{\overset{n}{\sum}}\frac{\ln^{(k)}(1)}{k!}t^{k}=t-\frac{t^{2}}{2}+\frac{t^{3}}{3}-\frac{t^{4}}{4}+\dots+\left(-1\right)^{n-1}\cdot\frac{t^{n}}{n}
\]
}%
\L{%
\[
\therefore\quad P_{n,\ln\left(1+x\right),0}(x^{2})=x^{2}-\frac{x^{4}}{2}+\frac{x^{6}}{3}-\dots
\]
}%
\L{%
\[
P_{n,\sin(x),0}(x)=\underset{k=0}{\overset{n}{\sum}}\frac{\sin^{(k)}(0)}{k!}x^{k}=x-\frac{x^{3}}{3!}+\frac{x^{5}}{5!}+\dots
\]
}%
\L{%
\[
\therefore\quad x\cdot P_{n,\sin(x),0}(x)=x^{2}-\frac{x^{4}}{3!}+\frac{x^{6}}{5!}+\dots
\]
}%

נפתח את הפונ׳ הראשונה עד {\beginL 2\endL} ואת הפונ' השנייה עד {\beginL 3\endL}.
נקבל\L{%
\[
\frac{\ln(1+x^{2})-x\sin(x)}{x^{n}}=\frac{x^{2}-\frac{x^{4}}{2}+R_{2,\ln\left(1+x\right),0}\left(x^{2}\right)-\left(x^{2}-\frac{x^{4}}{3!}+xR_{3,\sin\left(x\right),0}(x)\right)}{x^{n}}
\]
}%
\L{%
\[
=\frac{x^{2}-x^{2}-\frac{x^{4}}{2}+\frac{x^{4}}{3!}+R_{2,\ln\left(1+x\right),0}\left(x^{2}\right)-xR_{3,\sin\left(x\right),0}(x)}{x^{n}}
\]
}%
\L{%
\[
=\frac{-\frac{1}{3}x^{4}+R_{2,\ln\left(1+x\right),0}\left(x^{2}\right)-xR_{3,\sin\left(x\right),0}(x)}{x^{n}}
\]
}%

ננחש \L{$n=4$} ונקבל\L{%
\[
\lim_{x\to0}\frac{\ln(1+x^{2})-x\sin(x)}{x^{n}}=\lim_{x\to0}\frac{-\frac{1}{3}x^{4}+R_{2,\ln\left(1+x\right),0}\left(x^{2}\right)-xR_{3,\sin\left(x\right),0}(x)}{x^{n}}
\]
}%
\L{%
\[
\underset{(*)}{=}\underset{=-\frac{1}{3}}{\underbrace{\lim_{x\to0}\frac{-\frac{1}{3}x^{4}}{x^{4}}}}+\underset{=0}{\underbrace{\lim_{x\to0}\frac{R_{2,\ln\left(1+x\right),0}\left(x^{2}\right)}{\left(x^{2}\right)^{2}}}}-\underset{=\lim_{x\to0}\frac{R_{3,\sin\left(x\right),0}(x)}{x^{3}}=0}{\underbrace{\lim_{x\to0}\frac{xR_{3,\sin\left(x\right),0}(x)}{x^{4}}}}=-\frac{1}{3}
\]
}%

\L{$(*)$} מאריתמטיקה של גבולות, לפי כלל ההצבה של גבולות )\L{$x^{2}$}
לא מתאפס בסביבה מנוקבת של אפס( ולפי המשפט על התאפסות השארית של פולינום
טיילור כאשר מעלת המכנה שווה לדרגת הפיתוח של הפולינום.

עבור \L{$n\ge5$} נניח בשלילה שהגבול קיים במובן הצר. אז \L{%
\[
\lim_{x\to0}\frac{\ln(1+x^{2})-x\sin(x)}{x^{n}}=\underset{\in\mathbb{R}}{L}=\lim_{x\to0}\frac{-\frac{1}{3}+\frac{R_{2,\ln\left(1+x\right),0}\left(x^{2}\right)}{x^{4}}-\frac{xR_{3,\sin\left(x\right),0}(x)}{x^{4}}}{x^{n-4}}
\]
}%

נסמן\L{%
\[
g(x)=-\frac{1}{3}+\frac{R_{2,\ln\left(1+x\right),0}\left(x^{2}\right)}{x^{4}}-\frac{xR_{3,\sin\left(x\right),0}(x)}{x^{4}}
\]
}%

אז מצד אחד\L{%
\[
\lim_{x\to0}g(x)=\lim_{x\to0}\left(x^{n-4}\cdot\frac{g(x)}{x^{n-4}}\right)=\lim_{x\to0}\left(x^{n-4}\cdot\frac{\ln(1+x^{2})-x\sin(x)}{x^{n}}\right)\underset{\left(**\right)}{=}0
\]
}%

\L{$\left(**\right)$} לפי משפט חסומה כפול אפסה )פונ' המתכנסת במובן
הצר בנקודה חסומה בסביבה מנוקבת של הנק'(.

ומצד שני\L{%
\[
\lim_{x\to0}g(x)=-\frac{1}{3}
\]
}%

כפי שהראנו קודם.

קיבלנו סתירה, מכאן \L{$n=4$} זה ה-\L{$n$} המקסימלי כנדרש.

\newpage{}
\end{sol*}
\begin{problem}
מצאו \L{$q\in\mathbb{Q}$} כך ש-\L{$|\sqrt{e}-q|<10^{-3}$}.
\end{problem}

\begin{sol*}
צריך למצוא קירוב עם שגיאה של פחות מ- \L{$\frac{1}{10^{3}}$} עבור
\L{$e^{x}$} עם \L{$x=\frac{1}{2}$}.

נשתמש בנקודה \L{$x=0$} עבור הפיתוח.

לפי משפט טיילור עם השארית בצורת לגראנז' קיימת נקודה \L{$c\in\left(0,\frac{1}{2}\right)$}
כך ש-\L{%
\[
R_{n,e^{x},0}\left(\frac{1}{2}\right)=\frac{\left(e^{x}\right)^{\left(n+1\right)}(c)}{\left(n+1\right)!}0.5^{n+1}=\frac{e^{c}}{\left(n+1\right)!\cdot2^{n+1}}\underset{(*)}{\le}\frac{e^{\frac{1}{2}}}{\left(n+1\right)!\cdot2^{n+1}}\underset{\left(**\right)}{\le}\frac{2}{\left(n+1\right)!\cdot2^{n+1}}=\frac{1}{\left(n+1\right)!\cdot2^{n}}
\]
}%

\L{$(*)$} \L{$e^{x}$} מונוטונית עולה ו-\L{$c\in\left(0,\frac{1}{2}\right)$}. 

\L{$\left(**\right)$} הוכחנו בכיתה ש- \L{$e\le4$} אז \L{$\sqrt{e}\le2$}.

נבחר \L{$n=4$} ונקבל\L{%
\[
\left|e^{\frac{1}{2}}-P_{4,e^{x},0}\left(\frac{1}{2}\right)\right|=\left|R_{4,e^{x},0}\left(\frac{1}{2}\right)\right|\le\frac{1}{\left(4+1\right)!\cdot2^{4}}=\frac{1}{5\cdot4\cdot3\cdot2^{5}}=\frac{1}{60\cdot32}<\frac{1}{1000}
\]
}%
נפתח\L{%
\[
P_{4,e^{x},0}\left(x\right)=\underset{k=0}{\overset{4}{\sum}}\frac{\left(e^{x}\right)^{(k)}(0)}{k!}x^{k}=e^{0}+e^{0}x+\frac{1}{2}e^{0}x^{2}+\frac{1}{6}e^{0}x^{3}+\frac{1}{24}e^{0}x^{4}=1+x+\frac{x^{2}}{2}+\frac{x^{3}}{6}+\frac{x^{4}}{24}
\]
}%

לכן\L{%
\[
q=P_{4,e^{x},0}\left(\frac{1}{2}\right)=1+\frac{1}{2}+\frac{1}{8}+\frac{1}{48}+\frac{1}{384}\in\mathbb{Q}
\]
}%

מקיים את הנדרש וסיימנו.

\newpage{}
\end{sol*}
\begin{problem}
הוכיחו שהגבול \L{$\underset{x\to0}{\lim}\left(\cos(x)\right)^{1/x^{2}}$}
קיים, וחשבו את ערכו.
\end{problem}

\begin{sol*}
לכל \L{$x\in\left(-\frac{\pi}{2},\frac{\pi}{2}\right)\setminus\left\{ 0\right\} $}
מתקיים \L{$\cos(x)>0$} ולכן \L{%
\[
\underset{x\to0}{\lim}\left(\cos(x)\right)^{1/x^{2}}=\underset{x\to0}{\lim}e^{\ln\left(\cos(x)\right)\cdot\frac{1}{x^{2}}}=\underset{x\to0}{\lim}e^{\frac{\ln\left(\cos(x)\right)}{x^{2}}}
\]
}%

\L{$\underset{x\to0}{\lim}\cos(x)=1$} ו-\L{$\ln(x)$} רציפה לכן לפי
משפט ההרכבה של פונקציות רציפות \L{$\underset{x\to0}{\lim}\ln\left(\cos\left(x\right)\right)=\ln\left(1\right)=0$}.
בנוסף\L{%
\[
\underset{x\to0}{\lim}x^{2}=0,\quad\left(\ln\left(\cos\left(x\right)\right)\right)'=-\frac{\sin(x)}{\cos(x)},\quad\left(x^{2}\right)'=2x
\]
}%

ולכל \L{$x\in\left(-\frac{\pi}{2},\frac{\pi}{2}\right)\setminus\left\{ 0\right\} $}
מתקיים \L{$2x\ne0$} לכן לפי כלל לופיטל אם הגבול הבא קיים גם הגבול
הזה קיים והם שווים:\L{%
\[
\underset{x\to0}{\lim}\frac{\ln\left(\cos(x)\right)}{x^{2}}=\underset{x\to0}{\lim}\frac{-\frac{\sin(x)}{\cos(x)}}{2x}=\underset{x\to0}{\lim}\left(-\frac{\sin(x)}{2x\cos(x)}\right)=-\underset{x\to0}{\lim}\left(\frac{\sin(x)}{x}\cdot\frac{1}{2\cos(x)}\right)
\]
}%

הוכחנו בכיתה ש- \L{$\underset{x\to0}{\lim}\left(\frac{\sin(x)}{x}\right)=1$}
ומתקיים \L{$\underset{x\to0}{\lim}\left(\frac{1}{2\cos(x)}\right)=\frac{1}{2}$}
לכן \L{%
\[
-\underset{x\to0}{\lim}\left(\frac{\sin(x)}{x}\cdot\frac{1}{2\cos(x)}\right)=-\frac{1}{2}
\]
}%

\L{$e^{x}$} רציפה לכן לפי משפט ההרכבה של פונקציות רציפות\L{%
\[
\underset{x\to0}{\lim}\left(\cos(x)\right)^{1/x^{2}}=\underset{x\to0}{\lim}e^{\frac{\ln\left(\cos(x)\right)}{x^{2}}}=e^{\left(\underset{x\to0}{\lim}\frac{\ln\left(\cos(x)\right)}{x^{2}}\right)}=e^{-\frac{1}{2}}=\frac{1}{e^{\frac{1}{2}}}=\frac{1}{\sqrt{e}}
\]
}%
\end{sol*}

\end{document}
